% ========================= KPIs IMPLEMENTADOS ========================
\section{KPIs Implementados}
\label{sec:kpis}

Se definieron cinco indicadores, con la condición de que \textbf{cada base de datos aporte al menos uno}: tres se calculan desde Cassandra, uno desde MongoDB y uno desde Neo4j. La etapa 6 del DAG (\repo{src/kpis.py}{src/kpis.py}) los recalcula en cada corrida para el día de la hora procesada y los guarda en la tabla \texttt{kpis\_by\_date}, de donde los lee el \textit{dashboard}.

\begin{table}[H]
\centering
\caption{Resumen de los cinco indicadores y su valor para el lunes 12 de octubre de 2026}
\label{tab:kpis}
\vspace{0.2cm}
\footnotesize
\begin{tabular}{@{}clllr@{}}
\toprule
\textbf{KPI} & \textbf{Indicador} & \textbf{Base de datos} & \textbf{Fuente} & \textbf{Valor (12 oct)} \\
\midrule
1 & Eventos procesados por día       & Cassandra & \texttt{events\_by\_hour}          & 1,981,972 \\
2 & Crecimiento diario (flujo estable) & Cassandra & \texttt{events\_by\_hour}          & $-14.0$\,\% \\
3 & Actividad por usuario & Cassandra & \texttt{daily\_totals}            & 4.07 \\
4 & Tasa de fusión de PR & MongoDB & \texttt{payload.action}       & 93.8\,\% \\
5 & Repositorio más conectado        & Neo4j     & grado de \texttt{:Repo}       & 2,034 usuarios \\
\bottomrule
\end{tabular}
\end{table}

Los valores de ejemplo corresponden al 12 de octubre, un día con las 24 horas cargadas y con el flujo de referencias de Git completo. Para cada indicador se presenta su definición, su fórmula de cálculo y su interpretación, como exige el enunciado, seguidas de la ficha correspondiente del \textit{dashboard}.

% ------------------------------------------------------------------ KPI 1
\subsection{KPI 1: Número de Eventos Procesados por Día}

\paragraph{Definición.} Cantidad de eventos públicos de GitHub procesados por el \textit{pipeline} en el día $D$ (UTC).

\paragraph{Fórmula.}
\[
E_D \;=\; \sum_{h \,\in\, H_D} \; \sum_{t} \; n_{D,h,t}
\]
donde $n_{D,h,t}$ es el número de eventos del tipo $t$ en la hora $h$ del día $D$ (una fila de \texttt{events\_by\_hour}) y $H_D$ es el conjunto de horas cargadas de ese día.

\paragraph{Interpretación.} Mide el volumen de actividad ingerido y es la base de los demás indicadores. Solo es comparable entre días con las mismas horas cargadas y, como se mostró en la Sección~\ref{subsec:cobertura}, entre días con la misma cobertura de la fuente. Por eso la etapa 6 guarda junto al valor dos datos de contexto, que el \textit{dashboard} muestra en la tarjeta del indicador: los eventos del flujo estable y el número de horas del día que traen \textit{push}.

\paragraph{Valor.} $E_D = 1{,}981{,}972$ eventos el 12 de octubre, un 176\,\% por encima de la media de los 44 días completos (718,507). De ellos, 71,580 pertenecen al flujo estable, y el flujo de referencias de Git estuvo presente en las 24 horas.

\captura{dash-kpi1.png}{Ficha del KPI 1 en el \textit{dashboard}: definición, fórmula, cálculo para la fecha, lectura y evolución en el periodo.}

% ------------------------------------------------------------------ KPI 2
\subsection{KPI 2: Crecimiento Diario de Eventos}

\paragraph{Definición.} Variación porcentual, respecto al día anterior, del volumen de eventos del \textbf{flujo estable} (todos los tipos excepto \texttt{PushEvent}, \texttt{CreateEvent} y \texttt{DeleteEvent}), calculada sobre las horas cargadas en ambos días.

\paragraph{Fórmula.}
\[
G_D \;=\; \frac{S_D(H) - S_{D-1}(H)}{S_{D-1}(H)} \times 100,
\qquad S_D(H) \;=\; \sum_{h \,\in\, H} \; \sum_{t \,\notin\, I} n_{D,h,t},
\qquad H \;=\; H_D \cap H_{D-1}
\]
donde $I$ es el conjunto de los tres tipos intermitentes y $S_D(H)$ son los eventos del flujo estable del día $D$ sumando únicamente las horas de $H$. Si no hay datos del día anterior, el indicador no se calcula.

\paragraph{Interpretación.} $G_D > 0$ indica que la actividad de colaboración crece y $G_D < 0$ que disminuye. Al excluir los tipos intermitentes, un hueco de cobertura de la fuente no se lee como una caída ni como una subida de actividad.

\paragraph{Valor.} $G_D = (71{,}580 - 83{,}258) / 83{,}258 \times 100 = -14.0$\,\% frente al domingo 11 de octubre, sobre las mismas 24 horas. Con todos los tipos, ese mismo día habría marcado $+17.9$\,\%: el volumen total subió porque la fuente entregó más \textit{pushes}, mientras la actividad de colaboración bajaba.

\paragraph{Decisión de diseño.} El indicador pasó por dos correcciones. La primera versión comparaba totales diarios: un día a medio procesar aparecía con variaciones ficticias de más de 100 puntos, y el primer día del periodo mostraba un engañoso 0\,\%. Restringir la comparación a las horas comunes y dejar el valor nulo cuando no hay día previo eliminó ambos artefactos. La segunda corrección responde al hallazgo de la Sección~\ref{subsec:cobertura}. Con todos los tipos de evento, el indicador oscilaba entre $-92.6$\,\% y $+448.4$\,\%, y sus extremos coincidían con los días en que el flujo de referencias de Git aparecía o desaparecía de la fuente: el 7 de octubre marcaba $-89.2$\,\% cuando el flujo estable de ese día en realidad creció un 34.4\,\%, y el 25 de octubre marcaba $+169.5$\,\% con el flujo estable cayendo un 18.8\,\%. Calculado sobre el flujo estable, el indicador se mueve entre $-44.5$\,\% y $+57.4$\,\% en todo el periodo. Como Cassandra conserva los conteos por hora y tipo, el cambio no exigió recargar datos: los indicadores de los 44 días se recalcularon con \repo{scripts/recompute_kpis.py}{scripts/recompute\_kpis.py}.

\captura{dash-kpi2.png}{Ficha del KPI 2. Al calcularse sobre el flujo estable, la serie ya no reproduce los escalones de cobertura de la fuente.}

% ------------------------------------------------------------------ KPI 3
\subsection{KPI 3: Promedio de Actividad por Usuario}

\paragraph{Definición.} Promedio de eventos generados por cada usuario activo (actor distinto) en el día.

\paragraph{Fórmula.}
\[
A_D \;=\; \frac{E_D}{U_D}
\]
donde $U_D$ es el número de usuarios distintos del día (\texttt{daily\_totals.unique\_users}).

\paragraph{Interpretación.} Mide la intensidad de uso. Como la distribución de actividad es muy asimétrica, describe la intensidad agregada y no al usuario típico.

\paragraph{Valor.} $A_D = 1{,}981{,}972 / 486{,}477 = 4.07$ eventos por usuario. La cifra convive con una mediana de 1 evento por usuario: el 1\,\% de las cuentas más activas concentra el 43.4\,\% de los eventos del día, y las 4,019 cuentas automatizadas (sufijo \texttt{[bot]}) generan el 21.7\,\%. Una sola de ellas, \texttt{github-actions[bot]}, aporta 309,389 eventos, el 15.6\,\% del total.

\captura{dash-kpi3.png}{Ficha del KPI 3.}

% ------------------------------------------------------------------ KPI 4
\subsection{KPI 4: Tasa de Fusión de \textit{Pull Requests}}

\paragraph{Definición.} Porcentaje de los \textit{pull requests} cerrados en el día que fueron fusionados (aceptados).

\paragraph{Fórmula.}
\[
T_D \;=\; \frac{F_D}{F_D + N_D} \times 100
\]
donde $F_D$ es el número de eventos \texttt{PullRequestEvent} del día con \texttt{payload.action = "merged"} y $N_D$ el de aquellos con \texttt{payload.action = "closed"} (cerrados sin fusionar).

\paragraph{Interpretación.} Indica qué proporción del trabajo que se somete a revisión termina aceptada; es una medida de la salud del flujo de contribución.

\paragraph{Valor.} $T_D = 8{,}314 / (8{,}314 + 545) \times 100 = 93.8$\,\%. En todo el periodo se mantiene entre 88.8\,\% y 94.9\,\%, con media de 93.4\,\%.

\paragraph{Por qué MongoDB.} Es el único indicador que requiere un campo del \texttt{payload}. Cassandra solo conserva conteos por tipo de evento y no distingue un \textit{pull request} abierto de uno fusionado. La agregación se apoya en el índice \texttt{ix\_type\_date}:

\begin{lstlisting}[style=plano]
db.events.aggregate([
  { $match: { type: "PullRequestEvent",
              created_at: { $gte: ISODate("2026-10-12T00:00:00Z"),
                            $lt:  ISODate("2026-10-13T00:00:00Z") } } },
  { $group: { _id: "$payload.action", n: { $sum: 1 } } }
])
\end{lstlisting}

\captura{dash-kpi4.png}{Ficha del KPI 4. Es el indicador más estable del periodo porque depende del flujo de colaboración, cuya cobertura no varía.}

% ------------------------------------------------------------------ KPI 5
\subsection{KPI 5: Repositorio con Mayor Número de Conexiones}

\paragraph{Definición.} Repositorio con más usuarios distintos conectados en el grafo acumulado hasta el final del día $D$, es decir, el nodo \texttt{Repo} de mayor grado de entrada considerando solo las relaciones cuya primera interacción ocurrió hasta ese día.

\paragraph{Fórmula.}
\[
r^{*}_D \;=\; \arg\max_{r} \; \deg_D(r),
\qquad
\deg_D(r) \;=\; \bigl|\{\, u : (u) \rightarrow (r),\; t_{u,r} \le D \,\}\bigr|
\]
donde $\deg_D(r)$ cuenta los usuarios distintos con al menos una relación (\textit{push}, \textit{pull request}, \textit{issue}, estrella o \textit{fork}) hacia el repositorio $r$, y $t_{u,r}$ es la fecha de la primera interacción del usuario con el repositorio (propiedad \texttt{first\_seen} de la relación).

\paragraph{Interpretación.} Señala el proyecto que más comunidad concentra. Un grado anómalo, o un vecindario dominado por un solo tipo de relación, puede indicar popularidad puntual o actividad artificial en lugar de colaboración.

\paragraph{Valor.} $r^{*}_D = \texttt{morluto/rea}$, con $\deg_D(r^{*}_D) = 2{,}034$ usuarios al cierre del 12 de octubre. El indicador cambia con la fecha: el 1 de octubre el repositorio más conectado era \texttt{NousResearch/hermes-agent} (34 usuarios); del 2 al 6, \texttt{eternity4719/HowToLiveBetter} (199 al cierre del día 6); el 7 y el 8, \texttt{storytold/photocraft} (456); y desde el 9 de octubre, \texttt{morluto/rea}, que al corte reúne 6,843 usuarios. Como se analizó en la Sección~\ref{sec:neo4j}, el 82\,\% de esas conexiones son únicamente estrellas.

\begin{lstlisting}[style=plano]
MATCH (u:User)-[x]->(r:Repo)
WHERE x.first_seen < datetime('2026-10-13T00:00:00Z')
RETURN r.name AS repo, count(DISTINCT u) AS degree
ORDER BY degree DESC, repo LIMIT 1
\end{lstlisting}

\paragraph{Decisión de diseño y limitación.} En la primera versión el indicador no dependía de la fecha: la consulta recorría el grafo completo y devolvía el mismo valor para cualquier día. El filtro por \texttt{first\_seen}, propiedad que el modelo ya almacenaba, lo hizo comparable con los otros cuatro sin modificar el grafo. La limitación que subsiste es que \texttt{first\_seen} se fija al crear la relación con la fecha de la hora que se está cargando: si las horas no se cargan en orden cronológico (la carga histórica admite procesar varios días en paralelo), una relación puede quedar fechada después de su primera aparición real y el grado de las fechas intermedias, subestimado. El valor al corte no se ve afectado.

\captura{dash-kpi5.png}{Ficha del KPI 5: valor para la fecha elegida y, a la derecha, el ranking actual del grafo consultado en vivo.}
